\chapter{Estado del arte}

\noindent Este capítulo presenta los fundamentos teóricos y los trabajos previos que contextualizan este trabajo. En primer lugar, sitúa los mercados digitales de activos dentro de las economías de los videojuegos y diferencia los principales tipos de mercado y plataforma. Esta distinción permite entender que las posibilidades de intercambio, los precios observados y las reglas de cada operación dependen del entorno en el que se realiza.

Después, el capítulo introduce el aprendizaje supervisado como método para obtener estimaciones a partir de datos históricos. A continuación, presenta los fundamentos del aprendizaje por refuerzo en el caso de un único agente, incluyendo la interacción entre agente y entorno, los procesos de decisión de Markov, las políticas, las recompensas y el retorno acumulado. Estos conceptos se amplían después al aprendizaje por refuerzo multiagente, donde varios agentes actúan sobre un entorno compartido. Se revisan las formas de cooperación, las recompensas individuales y colectivas, y los enfoques centralizados y descentralizados.

Por último, se exponen los principios generales de la toma de decisiones en carteras y se revisan aplicaciones relacionadas de aprendizaje por refuerzo y aprendizaje por refuerzo multiagente. Esta revisión permite identificar las aportaciones, limitaciones y retos que fundamentan el planteamiento desarrollado en los capítulos posteriores.

\section{Mercados y plataformas de activos digitales}

Los bienes virtuales pueden tener utilidad dentro de un videojuego y, cuando las reglas lo permiten, circular entre personas usuarias a través de mecanismos de intercambio. La investigación sobre consumo de bienes virtuales muestra que la demanda no depende solo de una función dentro del juego. También intervienen la apariencia, la rareza, la identidad social y la percepción de valor del objeto \cite{lehdonvirta2009virtual,guo2009virtual}.

Por ejemplo, un atuendo para un avatar o una apariencia para un arma modifica su presentación visual sin cambiar las reglas de juego. Un objeto funcional puede dar acceso a una capacidad, una zona o un recurso dentro del juego, proporcionando ventaja al usuario. También existen objetos coleccionables cuyo interés depende de una serie limitada, su rareza o la demanda de la comunidad. Algunos bienes son consumibles, como cajas, sobres o mejoras, y desaparecen o se transforman al utilizarlos.

En los mercados de activos digitales hay dos formas de anunciar un intercambio:
\begin{itemize}
    \item Una persona usuaria que posee un activo digital y desea venderlo publica una \textbf{\glosref{oferta}{oferta}}. En esa publicación ofrece el activo digital e indica el precio de venta para que otra persona pueda comprarlo.
    \item Una persona usuaria que desea adquirir un activo digital publica una \textbf{\glosref{orden-compra}{orden de compra}}. En esa publicación ofrece una cantidad máxima de dinero por un activo digital concreto.
\end{itemize}

Cuando se completa el intercambio, la plataforma o el servicio de pago puede cobrar costes de transacción, como una comisión por la compra o la venta. Las reglas de cada plataforma determinan cómo se publican, consultan y aceptan estas publicaciones.

Para estudiar una posible compraventa, el precio de compra y el precio de venta deben expresarse en una moneda común. Las comisiones y la conversión de moneda modifican el resultado económico de la operación, por lo que una diferencia entre dos precios publicados no basta para determinar su beneficio neto. El volumen de intercambios indica cuántas unidades de un activo digital se han vendido durante un periodo y aporta información sobre la actividad observada para ese activo.

Algunas plataformas aplican un periodo de bloqueo de intercambio o \textit{trade hold}, durante el cual un activo digital recién adquirido no puede venderse o transferirse. Mientras el bloqueo permanece activo, el importe empleado en la compra constituye capital bloqueado y no puede reutilizarse en otra operación. Estas condiciones afectan al tiempo necesario para completar una compraventa y al saldo disponible de la cartera.

Los mercados de activos digitales pueden clasificarse en centralizados y descentralizados según quién administra la cuenta, los activos digitales asociados a ella y las reglas de intercambio:
\begin{itemize}
    \item \textbf{\glosref{mercado-centralizado}{Mercados centralizados}:} una empresa o plataforma administra las cuentas, los activos digitales o las reglas de intercambio. Dentro de esta categoría se encuentran:
    \begin{itemize}
        \item \textbf{\glosref{mercado-cerrado}{Mercados cerrados}:} los activos digitales se usan dentro del videojuego, pero no se intercambian entre personas usuarias. Por tanto, no constituyen un \textbf{\glosref{mercado-secundario}{mercado secundario}} (un mercado en el que las personas usuarias se revenden activos digitales entre sí) ni existe un precio de reventa observable.
        \item \textbf{\glosref{mercado-integrado}{Mercados integrados}:} el mercado está incluido dentro de la plataforma que administra la cuenta y los activos digitales asociados a ella. Esa plataforma establece las reglas para publicar ofertas y órdenes de compra.
        \item \textbf{\glosref{plataforma-externa}{Plataformas externas}:} proporcionan un espacio de intercambio distinto de la plataforma donde se gestionan originalmente la cuenta y los activos digitales. Aplican sus propias reglas, precios y mecanismos de pago.
    \end{itemize}
    \item \textbf{\glosref{mercado-descentralizado}{Mercados descentralizados}:} la gestión del activo y las reglas de intercambio no dependen de una plataforma central. Habitualmente emplean una infraestructura distribuida para registrar la propiedad y las operaciones.
\end{itemize}

Los mercados integrados están presentes en distintos ecosistemas de videojuego y distribución. La Tabla~\ref{tab:mercados-integrados} presenta algunos ejemplos.

\begin{table}[H]
\centering
\small
\begin{tabularx}{\textwidth}{@{}p{0.30\textwidth}p{0.23\textwidth}X@{}}
\toprule
\textbf{Mercado} & \textbf{Ecosistema} & \textbf{Activos habituales} \\
\midrule
Steam Community Market & Steam & Activos cosméticos, cajas y pegatinas \\
Roblox Marketplace & Roblox & Prendas digitales, accesorios y animaciones \\
Auction House & World of Warcraft & Equipo, materiales y moneda \\
EVE Online Market & EVE Online & Naves, módulos y recursos \\
\bottomrule
\end{tabularx}
\caption{Ejemplos de mercados integrados.}
\label{tab:mercados-integrados}
\end{table}

Aunque sus activos digitales y reglas son diferentes, estos mercados gestionan los activos digitales asociados a las cuentas y el mecanismo de intercambio dentro del mismo ecosistema \cite{steamMarketFAQ,robloxMarketplace,blizzardAuctionHouse,eveMarketOrders}.

Las plataformas externas ofrecen espacios de intercambio independientes para activos digitales transferibles. La Tabla~\ref{tab:plataformas-externas} presenta algunos ejemplos.

\begin{table}[H]
\centering
\small
\begin{tabularx}{\textwidth}{@{}p{0.20\textwidth}p{0.40\textwidth}X@{}}
\toprule
\textbf{Plataforma} & \textbf{Especialización} & \textbf{Modalidad} \\
\midrule
BUFF & Activos digitales transferibles & Compra y venta entre usuarios \\
Skinport & Activos cosméticos de videojuegos & Mercado con precios publicados \\
BitSkins & Activos digitales transferibles & Mercado de terceros \\
DMarket & Activos digitales de videojuegos & Compra y venta entre usuarios \\
\bottomrule
\end{tabularx}
\caption{Ejemplos de plataformas externas.}
\label{tab:plataformas-externas}
\end{table}

Las dos tablas no deben interpretarse como pares equivalentes, ya que cada una reúne ejemplos independientes de una configuración de intercambio \cite{buffMarketplace,skinportItemsApi,bitskinsMarketplace,dmarketMarketplace}.
\section{Aprendizaje supervisado y no supervisado}

El aprendizaje supervisado es un método que aprende a partir de ejemplos en los que se conoce tanto la información de entrada como el resultado que debe estimarse. Cada ejemplo se representa mediante una entrada $x$ y una etiqueta $y$, que es el resultado conocido asociado a esa entrada. Durante el entrenamiento, el modelo ajusta sus parámetros para aprender la relación entre ambas variables.

Una vez entrenado, el modelo aplica la relación aprendida a nuevas entradas y genera una estimación:

\begin{equation}
\hat{y}=f_{\theta}(x),
\label{eq:prediccion-supervisada}
\end{equation}

donde $x$ representa la información de entrada, $f_{\theta}$ el modelo con parámetros $\theta$ y $\hat{y}$ el resultado estimado. Cuando el resultado pertenece a una categoría, el problema se denomina clasificación. Cuando el resultado es un valor numérico, se denomina regresión. En un problema de clasificación binaria, el modelo también puede estimar la probabilidad de que una entrada pertenezca a la clase positiva:

\begin{equation}
\hat{p}=\mathbb{P}(Y=1\mid x).
\label{eq:probabilidad-supervisada}
\end{equation}

En la Ecuación~\ref{eq:probabilidad-supervisada}, $Y$ representa el resultado que se desea clasificar, $Y=1$ representa la clase positiva y $\hat{p}$ es la probabilidad estimada de pertenecer a esa clase. El entrenamiento consiste en seleccionar los parámetros que reducen el error entre las estimaciones y los resultados conocidos del conjunto de ejemplos. Si se dispone de $n$ ejemplos, este ajuste puede expresarse mediante una función de pérdida $\mathcal{L}$:

\begin{equation}
\theta^\star=\arg\min_{\theta}\frac{1}{n}\sum_{i=1}^{n}\mathcal{L}\bigl(y_i,f_{\theta}(x_i)\bigr).
\label{eq:perdida-supervisada}
\end{equation}

En la Ecuación~\ref{eq:perdida-supervisada}, $n$ es el número de ejemplos, $x_i$ es la entrada del ejemplo $i$, $y_i$ su resultado conocido y $\mathcal{L}$ mide el error entre ese resultado y la estimación del modelo. El valor $\theta^\star$ representa los parámetros que minimizan el error medio. La ecuación no fija un único modelo ni una única forma de medir el error. La función de pérdida se elige según la tarea. Por ejemplo, puede penalizar la diferencia entre un valor estimado y el valor observado en regresión, o penalizar una clasificación incorrecta cuando las etiquetas representan categorías.

El aprendizaje no supervisado parte de datos de entrada sin una etiqueta o resultado conocido asociado. Su objetivo es descubrir estructuras presentes en los datos, como grupos de observaciones con características parecidas, relaciones entre variables o comportamientos poco habituales. Por ello, puede utilizarse para explorar un conjunto de datos, agrupar elementos similares o detectar observaciones anómalas.

La diferencia esencial entre ambos enfoques se encuentra en la información disponible durante el entrenamiento. El aprendizaje supervisado aprende a relacionar una entrada con un resultado conocido, mientras que el aprendizaje no supervisado describe la estructura de los datos sin partir de un resultado objetivo. Ambos enfoques pueden utilizarse de forma complementaria, por ejemplo, cuando los grupos o patrones identificados mediante aprendizaje no supervisado se incorporan como información adicional a un modelo supervisado. Una estimación o una probabilidad no constituye por sí sola una decisión, ya que no incorpora automáticamente las restricciones ni las consecuencias posteriores de una acción.

\section{Aprendizaje por refuerzo}

El aprendizaje por refuerzo (\acroref{RL}{RL}, por sus siglas en inglés de \emph{Reinforcement Learning}) estudia problemas de decisión secuencial en los que un agente aprende mediante su interacción con un entorno \cite{sutton2018reinforcement}. En cada instante $t$, el agente recibe la información disponible, denominada estado $s_t$, selecciona una acción $a_t$ y recibe una recompensa $r_{t+1}$. Como consecuencia de esa acción, el entorno pasa a un nuevo estado $s_{t+1}$. A diferencia del aprendizaje supervisado, donde un modelo aprende a reproducir etiquetas históricas, en RL el agente aprende qué acciones conviene seleccionar a lo largo de una secuencia de decisiones.

Esta interacción puede representarse mediante un proceso de decisión de Markov, o \acroref{MDP}{MDP}, definido por el estado, las acciones, la transición, la recompensa y el factor de descuento:

\begin{equation}
\mathcal{M}=(\mathcal{S},\mathcal{A},P,R,\gamma),
\qquad 0\leq\gamma<1.
\label{eq:mdp-general}
\end{equation}

En la Ecuación~\ref{eq:mdp-general}, $\mathcal{M}$ representa el proceso de decisión de Markov, $\mathcal{S}$ es el conjunto de estados posibles y $\mathcal{A}$ el conjunto de acciones disponibles. La función de transición $P(s'\mid s,a)$ expresa la probabilidad de pasar del estado $s$ al estado $s'$ después de ejecutar la acción $a$. La función $R(s,a,s')$ asigna la recompensa obtenida en esa transición. El factor $\gamma$ determina cuánto peso reciben las recompensas futuras frente a las inmediatas.

El comportamiento que aprende el agente se denomina política y se representa por $\pi$. En una política estocástica, $\pi(a\mid s)$ expresa la probabilidad de seleccionar la acción $a$ cuando el agente se encuentra en el estado $s$:

\begin{equation}
\pi(a\mid s)=\mathbb{P}(A_t=a\mid S_t=s).
\label{eq:politica-rl}
\end{equation}

La Figura~\ref{fig:ciclo-rl} representa el ciclo básico de interacción. El estado del entorno se entrega al agente, el agente selecciona una acción y el entorno actualiza el estado y proporciona la recompensa correspondiente.

\begin{figure}[H]
  \centering
  \begin{tikzpicture}[
    node distance=1.7cm,
    bloque/.style={
      draw=tfmAccent,
      thick,
      rounded corners=3pt,
      fill=tfmSoft,
      minimum width=4.3cm,
      minimum height=1.15cm,
      align=center
    },
    flujo/.style={-{Latex[length=2.4mm]}, thick, draw=tfmAccent},
    retorno/.style={-{Latex[length=2.4mm]}, thick, draw=tfmMuted}
  ]
    \node[bloque] (estado) {Estado del entorno\\{\small $s_t$}};
    \node[bloque, below=of estado] (agente) {Agente\\{\small Política $\pi$}};
    \node[bloque, below=2.1cm of agente] (entorno) {Entorno\\{\small $P(s_{t+1}\mid s_t,a_t)$ y $R(s_t,a_t,s_{t+1})$}};

    \draw[flujo] (estado) -- node[right, font=\small] {estado $s_t$} (agente);
    \draw[flujo] (agente.south west) -- node[left, font=\small] {acción $a_t$} (entorno.north west);
    \draw[retorno] (entorno.north east) -- node[right, font=\small] {recompensa $r_{t+1}$} (agente.south east);
    \draw[retorno] (entorno.west) .. controls +(-4.0,0) and +(-4.0,0) .. node[left, font=\small, align=center] {nuevo estado\\$s_{t+1}$} (estado.west);
  \end{tikzpicture}
  \caption{Ciclo de interacción entre un agente y un entorno en aprendizaje por refuerzo.}
  \label{fig:ciclo-rl}
\end{figure}

El retorno mide la recompensa acumulada durante un episodio. Si $T$ representa el último instante del episodio, el retorno desde el instante $t$ se define como:

\begin{equation}
G_t^{\pi}=\sum_{k=0}^{T-t-1}\gamma^k r_{t+k+1}.
\label{eq:retorno-rl}
\end{equation}

En la Ecuación~\ref{eq:retorno-rl}, $r_{t+k+1}$ es la recompensa recibida después de la acción tomada en cada instante futuro y $\gamma^k$ es el descuento aplicado a esa recompensa. Por tanto, una recompensa inmediata tiene más peso que una recompensa situada varios instantes después. El objetivo de RL es aprender la política que maximiza el retorno esperado desde el estado inicial:

\begin{equation}
\pi^\star=\arg\max_{\pi}\mathbb{E}_{\pi}\left[G_0^{\pi}\right].
\label{eq:objetivo-rl-general}
\end{equation}

En la Ecuación~\ref{eq:objetivo-rl-general}, $\pi^\star$ es la política buscada y $\mathbb{E}_{\pi}$ representa el valor medio del retorno al seguir la política $\pi$. La ecuación no busca únicamente la acción con una recompensa inmediata mayor. Busca una política que tenga en cuenta las consecuencias de las decisiones posteriores. La propiedad de Markov no implica que el agente conozca el futuro, sino que el estado reúne la información necesaria para decidir en el instante actual.

Durante el aprendizaje, el agente debe equilibrar exploración y explotación. La exploración consiste en probar acciones cuya utilidad todavía es incierta para conocer sus consecuencias. La explotación consiste en seleccionar las acciones que, según lo aprendido hasta ese momento, tienen un retorno esperado mayor. Explorar en exceso puede acumular decisiones poco útiles, mientras que explotar demasiado pronto puede impedir descubrir alternativas mejores.

\section{Aprendizaje por refuerzo multiagente}

Un sistema multiagente reúne dos o más agentes que actúan sobre un mismo entorno. Estos agentes pueden colaborar, competir o tomar decisiones independientes. En un sistema cooperativo comparten un objetivo global, aunque cada agente disponga de información y responsabilidades distintas \cite{wooldridge2009multiagent,busoniu2008comprehensive}. El aprendizaje por refuerzo multiagente, o \acroref{MARL}{MARL}, amplía la formulación de \acroref{RL}{RL} para representar estos sistemas.

Un modelo genérico de MARL puede expresarse como una extensión del proceso de decisión de Markov:

\begin{equation}
\mathcal{M}=\left(\mathcal{N},\mathcal{S},\{\mathcal{A}^{i}\}_{i\in\mathcal{N}},P,\{R^{i}\}_{i\in\mathcal{N}},\gamma\right).
\label{eq:marl-general}
\end{equation}

En la Ecuación~\ref{eq:marl-general}, $\mathcal{N}$ es el conjunto de agentes, $\mathcal{S}$ es el conjunto de estados globales y $\mathcal{A}^{i}$ es el conjunto de acciones disponibles para el agente $i$. La función $P$ describe la transición del entorno después de la acción conjunta de los agentes. Cada función $R^{i}$ calcula la recompensa que recibe el agente $i$ y $\gamma$ es el factor de descuento. En cada instante, los agentes generan una acción conjunta:

\begin{equation}
\mathbf{a}_{t}=\left(a_t^1,\ldots,a_t^N\right),
\qquad
r_{t+1}^{i}=R^{i}\left(s_t,\mathbf{a}_{t},s_{t+1}\right).
\label{eq:acciones-recompensas-marl}
\end{equation}

En la Ecuación~\ref{eq:acciones-recompensas-marl}, $N$ es el número de agentes, $\mathbf{a}_{t}$ reúne las acciones tomadas en el instante $t$ y $r_{t+1}^{i}$ es la recompensa que recibe el agente $i$ tras la transición al estado $s_{t+1}$. Esta formulación permite distinguir cómo se reparten los incentivos entre los agentes.

En un problema plenamente cooperativo todos los agentes reciben la misma recompensa:

\begin{equation}
r_{t}^{1}=r_{t}^{2}=\cdots=r_{t}^{N}=r_t.
\label{eq:recompensa-compartida}
\end{equation}

En la Ecuación~\ref{eq:recompensa-compartida}, $r_t$ representa la recompensa común. Esta estructura alinea explícitamente a todos los agentes con el mismo resultado. En un problema cooperativo, en cambio, cada agente puede recibir una recompensa individual distinta. Estas recompensas no tienen por qué coincidir, pero deben diseñarse para que los comportamientos que incentivan contribuyan al objetivo compartido. La distinción es relevante porque una recompensa común favorece la cooperación directa, mientras que las recompensas individuales permiten reflejar responsabilidades diferentes y pueden dificultar la coordinación si no están bien alineadas.

MARL introduce dificultades que no aparecen con un único agente. Cada agente aprende y modifica su política mientras los demás también modifican las suyas. Por ello, desde el punto de vista de un agente, las consecuencias de una acción pueden cambiar aunque el entorno físico permanezca igual. Este fenómeno se denomina no estacionariedad. Además, un agente puede decidir a partir de una observación local que no contiene toda la información del estado global, lo que dificulta atribuir a cada decisión la parte del resultado que le corresponde.

La información disponible puede repartirse de forma distinta durante el entrenamiento y la ejecución. El entrenamiento es la fase en la que las políticas se ajustan a partir de la experiencia obtenida en el entorno. La ejecución es la fase posterior en la que las políticas ya entrenadas seleccionan acciones. En el entrenamiento centralizado y ejecución centralizada, conocido como \acroref{CTCE}{CTCE}, una entidad central recibe el estado global y selecciona la acción conjunta tanto durante el entrenamiento como durante la ejecución. Esta configuración facilita la coordinación porque concentra toda la información en una única decisión, pero exige que el estado global y la entidad central estén disponibles cuando el sistema actúa. Por ejemplo, un coche autónomo puede utilizar un controlador central que recibe toda la información disponible y decide la aceleración, el frenado y la dirección tanto durante el entrenamiento como durante la conducción.

El paradigma de entrenamiento centralizado y ejecución descentralizada, conocido como \acroref{CTDE}{CTDE}, separa ambas funciones. Durante el entrenamiento, un modelo evaluador puede utilizar el estado global y las acciones de todos los agentes para estimar la calidad esperada de su coordinación. Durante la ejecución, cada agente selecciona su acción a partir de su propia observación local. Por ejemplo, varios coches autónomos pueden entrenarse con un modelo evaluador que conoce la posición, la velocidad y las acciones de todos ellos. Cuando circulan, cada coche decide utilizando únicamente sus propios sensores. Este enfoque permite entrenar con información adicional sin exigir que esa información esté disponible cuando el sistema toma decisiones \cite{lowe2017maddpg,yu2022mappo}. El número de coches no determina el paradigma. Varios coches controlados por una única entidad central seguirían un enfoque CTCE. Por tanto, CTCE y CTDE representan dos formas distintas de distribuir la información y la responsabilidad de decisión entre los agentes.

\begin{figure}[H]
  \centering
  \begin{tikzpicture}[
    node distance=1.75cm and 1.35cm,
    bloque/.style={
      draw=tfmAccent,
      thick,
      rounded corners=3pt,
      fill=tfmSoft,
      minimum width=3cm,
      minimum height=1.05cm,
      align=center
    },
    evaluador/.style={
      bloque,
      draw=tfmMuted,
      fill=white,
      minimum width=4.8cm
    },
    flujo/.style={-{Latex[length=2.1mm]}, thick, draw=tfmAccent},
    entrenamiento/.style={-{Latex[length=2.1mm]}, dashed, thick, draw=tfmMuted}
  ]
    \node[bloque, minimum width=5.2cm] (estado) {Estado global y observaciones\\{\small $s_t$}};
    \node[bloque, below left=of estado] (agente1) {Agente 1\\{\small Observación local $o_t^1$}};
    \node[bloque, below=of estado] (agente2) {Agente $i$\\{\small Observación local $o_t^i$}};
    \node[bloque, below right=of estado] (agenteN) {Agente $N$\\{\small Observación local $o_t^N$}};
    \node[bloque, below=2.25cm of agente2, minimum width=5.2cm, minimum height=1.45cm] (entorno) {Entorno compartido\\{\small Acción conjunta $\mathbf{a}_t=(a_t^1,\ldots,a_t^N)$}\\{\small Recompensas $(r_{t+1}^1,\ldots,r_{t+1}^N)$}};
    \node[evaluador, below=1.65cm of entorno] (evaluador) {Modelo evaluador centralizado\\{\small Disponible solo durante el entrenamiento}};

    \draw[flujo] (estado) -- node[left, font=\scriptsize] {$o_t^1$} (agente1);
    \draw[flujo] (estado) -- node[right, font=\scriptsize] {$o_t^i$} (agente2);
    \draw[flujo] (estado) -- node[right, font=\scriptsize] {$o_t^N$} (agenteN);
    \draw[flujo] (agente1.south) -- node[left, font=\scriptsize] {$a_t^1$} (entorno.north west);
    \draw[flujo] (agente2.south) -- node[right, font=\scriptsize] {$a_t^i$} (entorno.north);
    \draw[flujo] (agenteN.south) -- node[right, font=\scriptsize] {$a_t^N$} (entorno.north east);
    \draw[entrenamiento] (entorno) -- node[right, font=\scriptsize, align=left] {estado global y\\acción conjunta} (evaluador);
  \end{tikzpicture}
  \caption{Interacción general en MARL con entrenamiento centralizado y ejecución descentralizada.}
  \label{fig:marl-ctde}
\end{figure}

La Figura~\ref{fig:marl-ctde} muestra una interacción general de este tipo. Las líneas continuas representan el intercambio necesario para ejecutar la política de cada agente. Las líneas discontinuas representan la información adicional que utiliza el modelo evaluador únicamente durante el entrenamiento.

\section{Gestión de carteras y riesgo}

La gestión de una cartera estudia cómo las decisiones sobre varios activos y el saldo disponible afectan al conjunto de recursos. No basta con evaluar cada activo de forma aislada. Una decisión de compra o venta de un activo puede parecer atractiva por sí sola y, sin embargo, puede aumentar demasiado la dependencia de la cartera respecto a un activo, una categoría o una plataforma.

La teoría moderna de carteras plantea que la rentabilidad y el riesgo deben analizarse conjuntamente \cite{markowitz1952portfolio}. La rentabilidad expresa el resultado obtenido en relación con los recursos utilizados. El riesgo recoge la incertidumbre y las posibles variaciones desfavorables que pueden afectar a ese resultado. Por ello, la gestión de carteras busca equilibrar ambos elementos en lugar de maximizar únicamente el resultado esperado de una decisión individual.

La asignación de recursos determina qué parte de la cartera se dedica a cada alternativa y qué parte permanece disponible. Una cartera concentrada destina una proporción elevada de sus recursos a pocos activos, categorías o plataformas. Esta concentración puede aumentar el resultado si los activos seleccionados evolucionan favorablemente, pero también aumenta el efecto de una variación desfavorable sobre el conjunto de la cartera.

La diversificación consiste en repartir los recursos entre alternativas distintas para reducir esa dependencia. No elimina el riesgo, pero evita que el resultado global dependa por completo de una sola alternativa. La exposición permite cuantificar esa dependencia al indicar qué parte de los recursos está vinculada a un activo, una categoría o una plataforma. Ambas nociones permiten comparar distintas distribuciones de una misma cartera.

La composición de la cartera no es necesariamente permanente. El reequilibrio consiste en revisar la distribución de los recursos y ajustarla cuando deja de responder al nivel de riesgo previsto. Esta revisión puede deberse a cambios en el valor de los activos, a la aparición de nuevas alternativas o a cambios en la información disponible. Al evaluar un reequilibrio deben considerarse los costes de transacción, ya que comprar o vender para modificar la cartera reduce el resultado obtenido.

También interviene el horizonte temporal, que determina durante cuánto tiempo se prevé mantener una decisión antes de revisarla. Un horizonte más amplio puede admitir variaciones de precio durante más tiempo, mientras que uno breve requiere que los recursos estén disponibles antes. 

Por tanto, la gestión de carteras combina la rentabilidad esperada, la concentración de los recursos, el plazo de las decisiones y los costes asociados a modificarlas.

\section{Aplicaciones y trabajos relacionados}

El aprendizaje automático se ha aplicado a problemas de análisis y decisión en los que los datos cambian con el tiempo. El aprendizaje supervisado permite estimar una variable a partir de observaciones históricas. El aprendizaje no supervisado permite descubrir agrupaciones o patrones en datos sin etiqueta. El aprendizaje por refuerzo permite estudiar decisiones encadenadas cuando una acción modifica las condiciones de las decisiones posteriores.

En la gestión automatizada de carteras, los métodos de aprendizaje por refuerzo se emplean para representar decisiones sucesivas de asignación, mantenimiento o modificación de recursos. La evaluación de estos métodos requiere conservar el orden temporal de los datos y reproducir las condiciones en las que se habría tomado cada decisión. También es necesario incorporar los costes de transacción, pues una estrategia que parece favorable antes de aplicarlos puede dejar de serlo después \cite{liu2021finrl}.

El aprendizaje por refuerzo multiagente se ha utilizado para estudiar situaciones en las que varios responsables de decisión comparten un mismo entorno. Estos responsables pueden cooperar, competir o combinar ambos comportamientos. La coordinación permite repartir una tarea compleja entre políticas distintas, aunque añade dificultades como la dependencia entre las decisiones de los agentes y la asignación del resultado colectivo a cada uno de ellos \cite{lowe2017maddpg,yu2022mappo}.

Estos trabajos muestran que la utilidad de un método no debe valorarse solo por el resultado que obtiene en los datos empleados para entrenarlo. También debe comprobarse su comportamiento con datos no utilizados durante el entrenamiento, su sensibilidad a cambios en las condiciones del entorno y la claridad de los supuestos de evaluación. Estos criterios permiten comparar métodos de aprendizaje y de decisión de forma reproducible.
